%%%%%%%%%%  scan idx 439  |  folio 426  %%%%%%%%%%
nous par exemple sur \(D_+(X_1)\). Alors \(p^*(E)\) est défini par le module
\[
 \mathbb Z[T_1,\ldots,T_d,X_2/X_1,\ldots,X_d/X_1]\otimes E,
\]
et \(\cO_P(1)\) par
\[
 \mathbb Z[T_1,\ldots,T_d,X_2/X_1,\ldots,X_d/X_1],
\]
le morphisme canonique envoyant \(e_i\) sur \(X_i/X_1\) pour tout \(i\). Alors \(F\) admet pour base les éléments \(e_i-e_1.X_i/X_1\) pour \(i>1\).
On a \(v(e_i-e_1.X_i/X_1)=T_i-T_1.X_i/X_1\) ; posons \(T'_i=T_i-T_1X_i/X_1\).
Comme la suite des \(T_i\) est régulière dans \(\mathbb Z[T_1,\ldots,T_d,X_2/X_1,\ldots,X_d/X_1]\), il en est de même de la suite \(T_1,T'_2,\ldots,T'_d\) qui s'en déduit par une matrice inversible, et donc aussi de la suite \(T'_2,\ldots,T'_d\) qui en est extraite (on est dans un cas noethérien). Par suite, \(v\) est un morphisme régulier.

Il faut donc montrer que les \(T'_i\) engendrent l'Idéal de \(j\). Dans \(\bigoplus_{n\ge0}J^n\), nous noterons, pour \(m\le p\), \(x^{(m)}\) l'élément \(x\) de \(J^p\), considéré comme étant de degré \(m\). Au-dessus de \(D_+(X_1)\), \(X'=\Spec\left(\left(\bigoplus_{n\ge0}J^n\right)_{(T_1^{(1)})}\right)\), et \(j\) est défini par le morphisme qui envoie \(X_i\) sur \(T_i^{(1)}\). Dans ce morphisme, l'image de \(T'_i\) est \(T_i^{(0)}-T_1^{(0)}.T_i^{(1)}/T_1^{(1)}=\bigl((T_iT_1)^{(1)}-(T_1T_i)^{(1)}\bigr)/T_1^{(1)}=0\). Donc les \(T'_i\) appartiennent à l'Idéal de \(j\). Comme tout élément de \(\mathbb Z[T_1,\ldots,T_d,X_2/X_1,\ldots,X_d/X_1]\) peut s'exprimer comme polynôme par rapport à \(T_1,T'_2,\ldots,T'_d,X_2/X_1,\ldots,X_d/X_1\), il suffit, pour montrer que les \(T'_i\) engendrent l'Idéal de \(j\), de montrer qu'aucun élément non nul de \(\mathbb Z[T_1,X_2/X_1,\ldots,X_d/X_1]\) n'appartient à l'Idéal de \(j\), ce qui revient à montrer que l'homomorphisme
\[
\mathbb Z[T_1,X_1,\ldots,X_d]\longrightarrow\bigoplus_{n\ge0}J^n
\]
est injectif, car \(T_1^{(1)}\) n'est pas diviseur de 0 dans \(\bigoplus_{n\ge0}J^n\). Si on gradue \(\mathbb Z[T_1,X_1,\ldots,X_d]\) en attribuant le degré \(0\) à \(T_1\), et le degré \(1\) aux \(X_i\), le noyau de l'homomorphisme considéré est un idéal gradué. Or il est clair que ce
%%%%%%%%%%  scan idx 440  |  folio 427  %%%%%%%%%%
noyau est contenu dans celui de l'homomorphisme composé :
\[
\mathbb Z[T_1,X_1,\ldots,X_d]\longrightarrow
\bigoplus_{n\ge0}J^n\longrightarrow A,
\]
le dernier étant défini en degré \(n\) par l'injection de \(J^n\) dans \(A\). Or le noyau de l'homomorphisme composé est l'idéal principal engendré par \((X_1-T_1)\), qui pour des raisons d'intégrité ne contient pas d'éléments homogènes non nuls. Donc le premier homomorphisme est injectif, ce qui achève la démonstration dans le cas considéré.

Revenons au cas général où \(A\) est quelconque ; on peut encore supposer \(E\) libre sur \(X\). Affectons l'indice \(0\) au cas particulier qui vient d'être traité ; alors d'après la partie i) de la proposition, la situation considérée est obtenue par le changement de base \(A_0\to A\) à partir de la situation étudiée précédemment, l'homomorphisme \(A_0\to A\) étant défini en prenant une base de \(E\), son image \(f_1,\ldots,f_d\) dans \(J\), et en envoyant \(T_i\) sur \(f_i\). Posons \(B_0=\left(\bigoplus_{n\ge0}J_0^n\right)_{(T_1^{(1)})}\), \(B=\left(\bigoplus_{n\ge0}J^n\right)_{(f_1^{(1)})}\) ; alors \(B=B_0\otimes_{A_0}A\) ; le complexe de Koszul défini par \(v:F\to\cO_P\) se déduit par changement de base du complexe de Koszul défini par \(v_0\) ; ce dernier étant une résolution projective de \(B_0\) (en se plaçant sur un ouvert convenable) les groupes d'homologie de \(K_\bullet(v)\) sont sur cet ouvert les \(\Tor_i^{A_0}(A,B_0)\). Ces derniers sont les composantes de degré \(0\) des \(\Tor_i^{A_0}\left(A,\left(\bigoplus_{n\ge0}J_0^n\right)_{T_1^{(1)}}\right)=\left(\bigoplus_{n\ge0}\Tor_i^{A_0}(A,J_0^n)\right)_{T_1^{(1)}}\) ; or les \(\Tor_i^{A_0}(A,J_0^n)\) sont nuls pour \(i\ge1\) d'après les relations 1.2, donc \(K_\bullet(v)\) est une résolution de \(B\), ce qui montre que \(j\) est régulière.

iii) La propriété étant locale sur \(X\), on peut encore supposer \(X\) affine et \(E\) libre de type fini. On reprend alors les notations de ii). On considère d'abord le cas standard de ii). L'anneau \(A_0\) est alors noethérien et de dimension cohomologique finie, et l'algèbre \(\cS(E_0)\)
%%%%%%%%%%  scan idx 441  |  folio 428  %%%%%%%%%%
également. On peut alors construire de proche en proche une résolution partielle graduée de \(\bigoplus_{n\ge0}J_0^n\), par des \(\cS(E_0)\)-modules gradués libres de type fini, de longueur supérieure à la dimension cohomologique de \(A_0\) ; le dernier noyau est alors un \(\cS(E_0)\)-module gradué, projectif de type fini en tant que \(\cS(E_0)\)-module non gradué, donc localement libre de type fini sur \(\Spec(A_0)\) en tant que \(\cS(E_0)\)-module gradué d'après VI 1.5.1. On obtient alors la résolution finie annoncée en remplaçant le dernier module de la résolution partielle par le noyau correspondant.

Si on tensorise la résolution ainsi construite par \(A\), on obtient alors un complexe gradué de \(\cS(E)\)-modules gradués localement libres de type fini, et il s'agit de voir que c'est une résolution de \(\bigoplus_{n\ge0}J^n\). Or le \(i\)-ième Module d'homologie du complexe ainsi obtenu est un \(\cS(E)\)-Module gradué dont la \(n\)-ième composante est \(\Tor_i^{A_0}(A,J_0^n)\), qui est nul pour \(i\ge1\) d'après (1.2.4). Par suite, d'après (1.2.3), on a construit ainsi une résolution du type voulu.

\textbf{Proposition 1.9.} Soient \(i:Y\to X\) une immersion fermée régulière, \(X'\) le \(X\)-schéma obtenu en faisant éclater \(Y\), et \(f:X'\to X\) le morphisme structural. Alors \(i\) et \(f\) sont des morphismes parfaits. (III 4.1.).

La propriété étant locale sur \(X\), on peut supposer donnés un \(\cO_X\)-Module localement libre de type fini \(E\), et un homomorphisme régulier \(u:E\to\cO_X\), ayant pour image l'idéal \(J\) de \(Y\). Le complexe de Koszul \(K_\bullet(u)\) est alors une résolution finie de \(\cO_Y\) par des \(\cO_X\)-Modules localement libres de type fini ; par conséquent, \(\cO_Y\) est un \(\cO_X\)-Module parfait. Il en résulte que toute immersion régulière est un morphisme parfait. D'après 1.8, \(f\) peut s'écrire comme le composé d'une immersion régulière \(X'\to\bfP(E)\) et du morphisme structural \(\bfP(E)\to X\) qui est lisse, donc parfait ; c'est donc le
%%%%%%%%%%  scan idx 442  |  folio 429  %%%%%%%%%%
composé de deux morphismes parfaits, et par suite un morphisme parfait.

\textbf{Proposition 1.10.} Soient \(X\), \(Y\) deux schémas lisses sur \(S\). Alors toute \(S\)-immersion de \(X\) dans \(Y\) est régulière. En particulier, toute section d'un morphisme lisse est une immersion régulière.

La proposition résulte immédiatement de EGA IV 17.12.1, puisqu'une immersion régulière au sens de EGA IV est régulière au sens de 1.4.

\section*{2. Calculs sur les immersions régulières}

\subsection*{2.1}
Soient \(A\) un anneau, \(B\) et \(C\) deux \(A\)-algèbres. On se propose de munir \(\bigoplus_{i\ge0}\Tor_i^A(B,C)\) d'une structure de \(A\)-algèbre graduée, donc de définir des homomorphismes :
\begin{equation}\tag{2.1.1}
 \Tor_i^A(B,C)\otimes_A\Tor_j^A(B,C)\longrightarrow \Tor_{i+j}^A(B,C).
\end{equation}
Soit \(P_\bullet\) une résolution de \(B\) par des \(A\)-modules projectifs; les \(\Tor_i^A(B,C)\) sont les groupes d'homologie du complexe \(P_\bullet\otimes_A C\). On sait définir des homomorphismes canoniques ([1], IV 6):
\begin{equation}\tag{2.1.2}
 H_i(P_\bullet\otimes_A C)\otimes_A H_j(P_\bullet\otimes_A C)
 \longrightarrow H_{i+j}\bigl((P_\bullet\otimes_A C)\otimes_A(P_\bullet\otimes_A C)\bigr).
\end{equation}
Par ailleurs, \(P_\bullet\otimes_A P_\bullet\) est un complexe de \(A\)-modules projectifs, à degrés \(\geq0\), et \(H_0(P_\bullet\otimes_A P_\bullet)=B\otimes_A B\); de plus, \(P_\bullet\) est une résolution de \(B\). Par suite, le morphisme d'augmentation \(B\otimes_A B\to B\) se prolonge en un morphisme de complexes : \(P_\bullet\otimes_A P_\bullet\to P_\bullet\) et définit donc, compte tenu de l'homomorphisme \(C\otimes_A C\to C\) déduit de la structure d'algèbre de \(C\), des homomorphismes :
\begin{equation}\tag{2.1.3}
 H_{i+j}\bigl((P_\bullet\otimes_A C)\otimes_A(P_\bullet\otimes_A C)\bigr)
 \longrightarrow H_{i+j}(P_\bullet\otimes_A C).
\end{equation}
%%%%%%%%%%  scan idx 443  |  folio 430  %%%%%%%%%%
En
composant (2.1.2) et (2.1.3), on obtient les homomorphismes (2.1.1). La structure d'algèbre définie par les homomorphismes (2.1.1) est associative (cf. [1] XI 2). Par ailleurs, le \(A\)-module \(H_i(P_\bullet\otimes_A C)\) est muni de façon standard d'une structure de \(B\otimes_A C\)-module en faisant opérer \(C\) de façon évidente, et \(B\) par fonctorialité; de même, les \(H_j\bigl((P_\bullet\otimes_A C)\otimes_A(P_\bullet\otimes_A C)\bigr)\) sont munis de deux structures de \(B\otimes_A C\)-modules. Il est clair que les homomorphismes (2.1.2) et (2.1.3) sont compatibles avec ces deux structures: pour (2.1.2), par fonctorialité (en se plaçant dans la catégorie des \(A\)-modules), et pour (2.1.3), en utilisant le résultat classique sur les morphismes d'un complexe acyclique. Nous avons donc muni \(\bigoplus_{i\ge0}\Tor_i^A(B,C)\) d'une structure de \(B\otimes_A C\)-algèbre associative.

\subsection*{2.2}
Supposons que \(B\) admette une résolution par des \(A\)-modules projectifs \(P_\bullet\), telle que \(P_\bullet=\bigoplus_{i\ge0}P_i\) soit une \(A\)-algèbre graduée, la différentielle satisfaisant:
\[
\begin{array}{ll}
\text{a)}& P_0\longrightarrow B\text{ est un homomorphisme d'anneaux~;}\\[0.4ex]
\text{b)}& \forall x\in P_i,\ \forall y\in P_j,\quad d(xy)=d(x)y+(-1)^i x\,d(y).
\end{array}
\]
\sgastandalonefootnotetext{x-tor-resolution}{On peut montrer qu'une telle résolution existe toujours.}
Alors les applications ``produit'' définissent des diagrammes commutatifs pour \(i,j\ge0\):
\begin{equation}\tag{2.2.1}
\begin{tikzcd}[column sep=large,row sep=large]
P_i\otimes_A P_j \arrow[r] \arrow[d] & (P_{i-1}\otimes_A P_j)\oplus(P_i\otimes_A P_{j-1}) \arrow[d]\\
P_{i+j} \arrow[r] & P_{i+j-1},
\end{tikzcd}
\end{equation}
et par suite un homomorphisme de complexes \(P_\bullet\otimes_A P_\bullet\to P_\bullet\) qui prolonge l'homomorphisme d'augmentation : \(B\otimes_A B\to B\). Ceci montre que la structure d'algèbre de \(\bigoplus_{i\ge0}\Tor_i^A(B,C)\) peut se calculer à partir de celle de \(P_\bullet\). En particulier, si \(P_\bullet\) est une \(A\)-algèbre alternée (i.e. telle que pour tout \(x\in P_\bullet\), on ait \(x^2=0\)),
%%%%%%%%%%  scan idx 444  |  folio 431  %%%%%%%%%%
